\scanpage{7}
\[
\begin{aligned}
0&\handwritten{\longrightarrow}\mathrm{Ker}\,v\xrightarrow{\handwritten{w}}\deleted{\unsure{(\mathrm{Coker}\,w)'}}\inserted{\handwritten{({}_p\mathrm{Tors}^{n}(X))_{ss}}}\handwritten{\longrightarrow}\\
&(E(Y)\otimes_W k)_{ss}\xrightarrow{\handwritten{\varphi}}
\deleted{\mathrm{Coker}\,v}\underset{\handwritten{\substack{\Vert\\\mathrm{Coker}\,v}}}{\handwritten{E(Y,\underline O_Y)_{ss}}}\handwritten{\longrightarrow}
\deleted{\unsure{(\mathrm{Ker}\,w)'\longrightarrow0}}\handwritten{0}\;,
\end{aligned}
\tag*{(8.5.1)}
\]
où on désigne par $v$ l'homomorphisme de restriction
\[
v\mathbin{:}H^{n-1}(X,\underline O_X)_{ss}\handwritten{\longrightarrow}H^{n-1}(Y,\underline O_Y)_{ss}\;,
\tag*{(8.5.2)}
\]
(en particulier $\mathrm{Coker}\,v$ est la ``cohomologie évanescente cohérente'' $E(Y,\underline O_Y)$), \deleted{et par $w$ l'homomorphisme}
\margin{\handwritten{\deleted{\unsure{$\xi=c_X$|$\xi=c_1X$}}}}
\[
\deleted{w\mathbin{:}(\mathrm{Tors}^{n-1}(X)\otimes_W k)_{ss}\xrightarrow{\handwritten{\unsure{\cup d\xi|u_{d\xi}}}}(\mathrm{Tors}^{n+1}(X)\otimes_W k)_{ss}}
\tag*{\deleted{(8.5.3)}}
\]
\margin{\handwritten{Katz fait \unsure{observer} que \unsure{$V$} \illegible{1} \unsure{est tjrs} \unsure{nul} \unsure{?} sur la \illegible{1}}}
\deleted{induit par l'homomorphisme $H^{n\deleted{\unsure{+1}}\handwritten{-1}}(X)\handwritten{\longrightarrow}H^{n+1}(X)$, \deleted{\unsure{cup-}}défini par le cup-produit avec la classe de polarisation \inserted{\handwritten{$d\xi$}} sur $X$.} Enfin, pour utiliser cette suite exacte (8.5.1), on tiendra compte du fait que
\[
E(Y)\overset{\rm dfn}{=}\mathrm{Coker}(H^{n-1}(X)\handwritten{\longrightarrow}H^{n-1}(Y))\quad\deleted{\text{est}}
\tag*{(8.5.4)}
\]
est un $W$-module libre \uline{sans torsion}; en particulier, si $k$ est un corps fini, le rang de $(E(Y)\otimes_W k)_{ss}$ est égal au nombre des valeurs propres de Frobenius géométrique relatif à $k$ opérant sur $\deleted{H}E^{n-1}(\deleted{\unsure{X}}Y)$ (la cohomologie cristalline évanescente de $Y$; ou ce qui revient au même, la cohomologie $\handwritten{l}$-adique pour n'importe quel $\handwritten{l\ne}p$ - sous réserve de l'algébricité de l'opération $\handwritten{\Lambda}$ sur lescohomologies cristallines et $\handwritten{l}$-adiques) \uline{qui sont des unités $p$-adiques}; les valeurs propres en question, réduites modulo $p$, sont égales aux valeurs\handwritten{\,}propres du frobenius $\deleted{\unsure{p}}\unsure{k|p}$-linéaire (relatif au corps de base fini $k$) opérant sur le terme $E(Y)\otimes_W k$ de (8.5.1), ou si on préfère, en termes de l'interprétation 8.4.2, aux valeurs propres d\deleted{u}e \inserted{l'inverse du} générateur de frobenius de $G=\mathrm{Gal}(\bar k/k)$ opérant sur le vectoriel de dimension finie sur $\underline{\rm F}_p$ correspondant.

Notons d'autre part l'égalité suivante déduite de (8.5.1):
\[
\begin{aligned}
\dim(E(Y)\otimes_W k)_{ss}={}&\dim E(Y,\underline O_Y)_{ss}\ \deleted{\unsure{-\dim\mathrm{Ker}\,w}}\\
&+(\dim\deleted{\mathrm{Coker}\,w}\inserted{\handwritten{({}_p\mathrm{Tors}^n(X))_{ss}}}-\dim\mathrm{Ker}\,v)\handwritten{,}\ \deleted{\unsure{+}}
\end{aligned}
\tag*{(8.5.5)}
\]
où le terme entre parenthèses du deuxième membre est $\handwritten{\geq}0$, de s\deleted{\unsure{o}}orte qu'\unsure{on}
% archive sheet number: 6
